% Generated by tools/paper/render.py from main_text.py; edit the source there, not this file.
\section{Discussion}

dinara-align was the fastest exact aligner compared on every dataset at unit costs, and at gap-affine costs on every dataset that any tool completed. This is consistent with choosing an algorithm for each pair: no single engine was fastest across these workloads, and since every engine is exact, choosing never forfeits the guarantee or changes the alignment. Three observations stand out. First, each technique paid where its engine does the work (Supplementary Table S2): on the ultra-long reads, switching off direct neighbour construction, which reduces work outside the inner loop, slowed alignment by 9 and 16\%, more than either change to the loop itself (at most 5\%), while the interleaved groups paid most on the widest bands and gathered extension where diagonal transition does most of the work. Second, on short reads the work that does not depend on the pair dominates: a 250~bp alignment at gap-affine costs once took 5.2~\textmu s through a fresh call and 2.4~\textmu s through a reused aligner (Supplementary Table S1), so a benchmark should call every tool as its own runner does. Third, thresholds that encode a cost model become outdated as the code changes, and measuring them again belongs in the evaluation of any aligner that selects its algorithm by thresholds. Gap-affine alignment of ultra-long reads remains open: no tool completed an ont-500k read within 30~s. Where memory matters more than time, Edlib, QuickEd and BiWFA (15 to 29 MiB on ultra-long reads) use less than dinara-align's 71 to 81 MiB. The measured-work rule, the regrouped recurrence and the per-lane certificate do not depend on this implementation and apply to other wavefront, bit-parallel and inter-sequence aligners. And because its alignments are exact and the same whichever engine computes them, dinara-align can serve as the reference against which heuristic aligners are measured, at the scale of current long reads.

\subsection{Limitations}

\begin{enumerate}
  \item The evaluation covers one x86 processor with AVX-512, and single-thread latency except for the batches; the AVX2 and NEON configurations are not evaluated here.
  \item The evaluated samples also guided development: several changes were motivated and checked on the Illumina, HiFi and ont-500k-genvar samples, and some thresholds, in particular the AVX-512 gates for inexact seeds, were set on uniform 100~kbp pairs like those of Fig.~\ref{fig:sweeps}a, so results on unseen data may be less favourable.
  \item The ultra-long and the largest simulated sets were sampled to four pairs. The three rounds were timed from one build, while on the Skylake-X running times varied by up to a third between builds that left the code involved unchanged (Supplementary Table S2), so margins of less than about 10\% may not hold across builds.
  \item The Illumina, HiFi and PromethION reads come from one region of one chromosome each and exclude reads that minimap2 leaves more than a tenth clipped, such as those spanning structural variants.
  \item The ablation switches off one technique at a time, so it does not separate their interactions.
\end{enumerate}

\subsection{Future work}

\begin{enumerate}
  \item A cheap upper bound on the distance. With the true distance supplied by an oracle, starting band doubling at it saves up to 27\% on 100~kbp pairs at 20\% errors, but starting 15\% above it loses up to 23\% on ultra-long reads, whereas capping the bound at an upper bound within 5\% of the distance saves up to 27\% on 1~Mbp pairs at 20\% (Supplementary Section S5).
  \item Reusing computation across band-doubling rounds, as in the incremental doubling of A*PA2, once its exactness for these tiles is proven.
  \item Faster construction of the seed heuristic, whose scalar lookups dominate on long divergent pairs, and a seed heuristic for gap-affine costs, at which no tool completed a 1~Mbp pair beyond 1\% errors.
  \item Measurements on further processors, and an evaluation of the GPU kernels the library provides for batches (Supplementary Section S1.5).
\end{enumerate}
